% !TeX root = ../build/main.tex

In this section we describe the overall flow of the \davinci voting protocol. An election is structured into five main phases: election setup, encryption key generation, voting period, tally decryption, and election finalization. Each phase specifies the interactions between the organizer, \keywardens, voters, sequencers, and the smart contracts that coordinate the election. %An overview of the complete process is shown in Fig.~\ref{fig:protocol-flow}.

\subsubsection{Election setup}
\label{sec:election-setup}

%Before any election can begin, sequencers must register in the \hlget{\electionregistry} smart contract. This registration is a one-time
%operation per sequencer and is independent of any particular election.
%; only registered
%sequencers are eligible to participate in the processing of votes.
%
For each new election, the organizer performs the following steps:

\begin{enumerate}
	
	\item Prepare the census data \censusdata\ and build the incremental Merkle tree
	$\censusTree$ from it. Publish the full tree to IPFS, and record the resulting URI as $\censusuri$.
	
	\item Compute the census commitment $\censusroot = \censusTree.\hlget{\MTRoot}()$.
	
	\item Define the election details: the start time (\starttime), the duration of the voting period (\duration), the census model $\censusorigin$
	(see Section~\ref{sec:vocdoni-protocol:census}), and the ballot rules $\ballotmode$
	(see Section~\ref{sec:ballot-protocol}).
	
	\item Upload the election metadata (questions, options, description, etc.)\ to IPFS and
	record the resulting URI as $\electionmetadata$.
	
	\item Determine the process identifier $\processid$, a 32-byte value derived
	deterministically from the organizer's Ethereum address, the account nonce, and a prefix
	encoding the chain identifier and the \hlget{\electionregistry} smart contract. Being deterministic, it can be
	computed in advance by the organizer, although it is ultimately recomputed and assigned by
	the contract upon registration.
	
	\item Obtain the encryption public key $\epk$ for this process. The key is produced by the
	key wardens through the \ac{DKG} protocol 
	run for the identifier $\processid$, so that no single party knows the corresponding secret key.
	
	\item Register the process on-chain by submitting a transaction to the
	\hlget{\electionregistry} smart contract with the election parameters:
	\begin{align*}
		\mathsf{tx} = \big(
		&\censusorigin,\ \censusroot,\ \censusuri,\ \electionmetadata, \\
		&\ballotmode,\ \epk,\ \starttime,\ \duration\big).
	\end{align*}
	Upon accepting this transaction, the \hlget{\electionregistry} contract initializes the reserved configuration leaves
	of the state tree $\stateTree$ as %
	$\processid$ at $\mathtt{0x0}$, $\ballotmode$ at $\mathtt{0x2}$, $\epk$ at $\mathtt{0x3}$,
	$\censusorigin$ at $\mathtt{0x6}$, and the results accumulator $\encresults$ at
	$\mathtt{0x4}$ initialized to the encryption of zero  (see Section~\ref{sec:vocdoni-protocol:state-tree}) and computes the initial state root
	\[
	\stateroot = \stateTree.\hlget{\MTRoot}().
	\]
\end{enumerate}
%
Once this transaction is accepted, the election is publicly registered on-chain and, at its
scheduled start time, the voting period can begin.
%
%Before any election can begin, key wardens must register in the \hlget{\dkgmanagement} smart contract. This registration
%is a one-time operation per key warden and is independent of any particular election; only
%registered parties are eligible to participate in the \ac{DKG} phase. 
%%
%For each new election, the organizer performs the following steps:
%
%\begin{enumerate}
%	
%	\item Prepare the census data \censusdata\ and build the \ac{IMT}
%	$\censusTree$ from it. Publish the full tree to IPFS, and record the resulting URI as $\censusuri$.
%	
%	\item Compute the census commitment $\censusroot = \censusTree.\hlget{\MTRoot}()$.
%	
%	\item Define the election details, including duration, the census model $\censusorigin$
%	(see Section~\ref{sec:vocdoni-protocol:census}), and the ballot rules $\ballotmode$
%	(see Section~\ref{sec:ballot-protocol}). 
%	
%	\item Upload the election metadata (questions, options, description, etc.)\ to IPFS to obtain {\color{blue}electionMetadata}.
%	
%	\item Compute a unique process identifier $\processid$, which is a 32-byte number derived
%	from the organizer's Ethereum address, nonce, and the chain ID.
%	
%	\item Initialize the state Merkle tree $\stateTree$ by inserting the election configuration
%	parameters into the reserved configuration namespace (see Section~\ref{sec:vocdoni-protocol:state-tree}):
%	$\processid$ at leaf $\mathtt{0x0}$, $\ballotmode$ at leaf $\mathtt{0x2}$,
%	$\censusorigin$ at leaf $\mathtt{0x6}$, and zero-valued accumulators $\addaccumulator$ and
%	$\substractaccumulator$ at leaves $\mathtt{0x4}$ and $\mathtt{0x5}$, respectively. {\color{blue}Compute the
%	initial state root -- is this done by the SC?}
%	\[
%	\stateroot = \stateTree.\hlget{\MTRoot}().
%	\]
%	
%	\item Set security parameters for the \ac{DKG} ceremony (e.g.\ timeout, minimum number
%	of sequencers) {\color{blue}DKG}.
%	
%	\item Submit all this information in a transaction to the \hlget{\electionregistry} smart
%	contract:
%	\begin{align*}
%		tx_{\mathit{org}} = \big(
%		&\processid,\ \censusroot,\ \censusorigin,\ \censusuri, \\
%		&\ballotmode,\ {\color{blue}electionMetadata},\ \stateroot,\ {\color{blue}dkgParameters}
%		\big).
%	\end{align*}
%	
%\end{enumerate}
%%
%Once this transaction is accepted, the election is publicly registered on-chain and the
%\ac{DKG} phase can begin.

\subsubsection{Encryption key generation}
\label{sec:vocdoni-protocol:keygen}

In this phase, the \keywardens\ collectively generate the public encryption key ($\epk$) that voters will use to encrypt their ballots. Since the process identifier $\processid$ is derived deterministically (see Section~\ref{sec:election-setup}), it is known in advance, so the \ac{DKG} ceremony for the election can be carried out before the process is registered on-chain.

Each \keywarden\ contributes to the ceremony and publishes a proof that its contribution is well-formed. The \hlget{\dkgmanagement} smart contract verifies these proofs, so that correct participation can be checked without revealing any secret. Once a threshold of valid contributions is collected, the joint public key $\epk$ is derived, while the corresponding secret key remains secret-shared among the wardens. In this way, no single party can decrypt on its own, and later tally decryption requires the collaboration of a threshold of wardens (see Section~\ref{sec:vocdoni-protocol:tally}). The resulting $\epk$ is the key that the organizer includes when registering the election (see Section~\ref{sec:election-setup}). The concrete \ac{DKG} construction is given in Appendix~\ref{sec:cryptographic-primitives:dkg}.

\subsubsection{Voting period}
\label{sec:vocdoni-protocol:votes}

In this phase, voters cast their ballots and send them to the sequencers, who collect and process them. \\

\noi \textbf{(a) Vote casting.} Each voter holds an Ethereum key pair $(\sk, \pk)$ with
associated address $\address$ (see Appendix~\ref{sec:cryptographic-primitives:signatures}).
To cast a ballot, a voter does the following:

\begin{enumerate}
	\item Retrieve the census membership proof $\censusproof$ to prove eligibility.
	\item Fetch the parameters $\epk$, $\ballotmode$, and $\processid$ from the
	\hlget{\electionregistry} contract.
	\item Select their ballot choice $\ballot$ following the rules of $\ballotmode$.
	\item Sample fresh randomness $k \in \F$ and encrypt the ballot using $\epk$:
	\begin{equation}\label{eq-ciphertext}
		\encryptedballot = \hlget{\Enc}_{\epk}(\ballot; k).
	\end{equation}
	\item Use their address $\address$ and the randomness $k$ to compute a unique vote
	identifier, mapped into the reserved vote-identifier namespace
	(see Section~\ref{sec:vocdoni-protocol:state-tree}):
	\begin{equation}\label{eq-voteid}
		\voteid = N + \bigl(\hlget{\Poseidon}(\processid, \address, k) \bmod N\bigr).
	\end{equation}
	\item Use the \ballotcircuit from Section~\ref{sec:vocdoni-protocol:circuits:ballot} to
	generate a \ac{ZK-SNARK} proof that Eqs.~\labelcref{eq-ciphertext,eq-voteid} are correctly
	computed and that $\ballot$ satisfies the ballot rules of $\ballotmode$:
	\begin{align*}
		\ballotproof = {} & \hlget{\PProve}(\hlget{\ballotcircuit}, \witness = (\ballot, k), \\
		& \publicinputs = (\processid, \ballotmode, \epk, \address, \weight,
		\encryptedballot, \voteid)).
	\end{align*}
	\item Sign the vote identifier with their Ethereum secret key $\sk$ to authenticate the vote:
	\[ \signature = \hlget{\SSign}_{\sk}(\voteid). \]
	\item Select a sequencer and submit the package
	\[ \vote = [\processid, \voteid, \encryptedballot, \censusproof, \ballotproof, \signature] \]
	to the chosen sequencer.
\end{enumerate}

\noi \textbf{(b) Vote batching.} Upon receiving votes, a sequencer does the following:

\begin{enumerate}
	\item Retrieve the election identifier $\processid$, the current state root $\stateroot$, and
	the census root $\censusroot$ from the \hlget{\electionregistry} smart contract, together with
	the associated data from the Ethereum data blobs.
	\item Upon receiving a vote
	\[ \vote_i = [\processid, \voteid_i, \encryptedballot_i, \censusproof_i, \ballotproof_i, \signature_i], \]
	recover the voter's public key $\pk_i = \hlget{\SPubKey}(\signature_i)$.
	\item For each $\vote_i$, use the \verifiercircuit from
	Section~\ref{sec:vocdoni-protocol:circuits:verifier} to generate a proof
	\begin{align*}
		\authenticationproof_i = {} & \hlget{\PProve}(\hlget{\verifiercircuit},
		\witness = (\pk_i, \signature_i, \ballotproof_i), \\
		& \publicinputs = (\processid, \ballotmode, \epk, \encryptedballot_i, \voteid_i,\\
		&\address_i, \weight_i)).
	\end{align*}
	This proof ensures that $\ballotproof_i$ and $\signature_i$ are valid and that $\address_i$
	is correctly derived from $\pk_i$. The census membership proof $\censusproof_i$ is not
	used here; it is retained for the state transition circuit
	(Section~\ref{sec:vocdoni-protocol:circuits:state-transition}), where eligibility is enforced.
	\item Batch a set of $n$ authentication proofs $\{\authenticationproof_i\}_{i=1}^{n}$ and verify
	them together using the \aggregationcircuit from
	Section~\ref{sec:vocdoni-protocol:circuits:aggregation}:
	\begin{align*}
		\aggregationproof = {} & \hlget{\PProve}(\hlget{\aggregationcircuit},
		\witness = (\{\authenticationproof_i\}_{i=1}^{n}), \\
		& \publicinputs = (\{\encryptedballot_i\}_{i=1}^{n}, \{\voteid_i\}_{i=1}^{n},
		\{\address_i\}_{i=1}^{n}, \processid, \\
		&\ballotmode, \epk)).
	\end{align*}
	This proof ensures that all $\authenticationproof_i$ are valid and share the same global
	parameters $\processid$, $\ballotmode$, and $\epk$.
	\item Prepare the state transition data to be published in an Ethereum blob: pack the votes and
	results into a 4096-cell blob $\blobdata$ and
	\begin{itemize}
		\item compute the data commitment $\blobcom = \hlget{\CCommit}(\blobdata)$,
		\item derive the evaluation point $\evalpoint = \hlget{\Poseidon}(\processid, \stateroot^{\old}, \blobcom)$,
		\item build the polynomial $P$ from $\blobdata$ and evaluate $y = P(\evalpoint)$,
		\item produce the opening proof $\openproof$ for the evaluation $(\evalpoint, y)$.
	\end{itemize}
	The consistency $y = P(\evalpoint)$ between the blob and the on-chain state is proven inside
	the state transition circuit.
	\item Generate the state transition proof with the \statetransitioncircuit from
	Section~\ref{sec:vocdoni-protocol:circuits:state-transition}:
	\begin{align*}
		\stateproof = {} & \hlget{\PProve}(\hlget{\statetransitioncircuit},
		\witness = (\aggregationproof, \{\censusproof_i\}_{i=1}^{n},\\
		&\{\encryptedballot_i\}_{i=1}^{n}, k, \openproof, \dots), \\
		& \publicinputs = (\censusroot, \stateroot^{\old}, \stateroot^{\new}, \blobcom, \dots)).
	\end{align*}
	The complete input list is given in
	Section~\ref{sec:vocdoni-protocol:circuits:state-transition}. This proof attests that the
	encrypted votes are correctly accumulated via ElGamal's homomorphic properties, that every
	processed vote corresponds to an eligible voter (via the census membership proofs), that vote
	overwrites are handled consistently, and the blob data matches the committed on-chain state.
	\item Submit the state transition proof $\stateproof$ and the new state root $\stateroot^{\new}$
	to the \hlget{\electionregistry} smart contract, while the full data is published in the
	Ethereum blobs.
\end{enumerate}

\noi Upon receiving the transaction, the smart contract checks that the submitted previous root
matches the current on-chain state root and that $\stateproof$ is valid before updating the state
root to $\stateroot^{\new}$. If the on-chain state root changed while the sequencer was preparing
the batch (because another sequencer committed first), only the state transition proof must be
regenerated, as noted in Section~\ref{sec:vocdoni-protocol:circuits:state-transition}. Sequencers
repeat this process until the voting deadline set by the organizer.

\subsubsection{Tally decryption}
\label{sec:vocdoni-protocol:tally}

After the voting period expires, the \hlget{\electionregistry} smart contract stops accepting new state updates, fixing the final state root and thus the aggregated encrypted results. A threshold of $t$ out of $n$ \keywardens\ then cooperate to decrypt this aggregate: rather than reconstructing the election secret key, each participating warden publishes a partial decryption of the encrypted results together with a proof that its contribution is correct with respect to its key share. Once $t$ valid partial decryptions are available, they are combined to recover the plaintext tally. Because only the aggregated results are decrypted and never the individual ballots, voter privacy is preserved throughout the election, and the secret key is never reconstructed by any single party. 
%
The resulting tally, together with a proof of correct decryption, is published on-chain during the finalization phase (see Section~\ref{sec:vocdoni-protocol:finalization}), where it is checked by the \resultscircuit\ and the \hlget{\resultsverifier} contract. %\Keywardens\ that fail to provide a valid partial decryption may be penalized; the incentive and slashing mechanisms are described in Section~\ref{sec:token}.
	

\subsubsection{Election finalization}
\label{sec:vocdoni-protocol:finalization}

In this final phase, the organizer (or a sequencer) publishes the election results, together with
a proof of their correct computation, on-chain.

\begin{enumerate}
	\item Obtain the plaintext results $\{\text{results}_j\}_{j=1}^{n_f}$ and the corresponding
	decryption proofs $\{\pi^{\mathrm{dec}}_j\}_{j=1}^{n_f}$ from the tally decryption phase
	(Section~\ref{sec:vocdoni-protocol:tally}), i.e.\ the combination of the \keywardens'
	partial decryptions of the aggregated encrypted results.
	\item Retrieve the encrypted results accumulator $\encresults$ and the encryption key $\epk$
	from the final state, together with their Merkle inclusion proofs against the final state
	root $\stateroot$.
	\item Use the \resultscircuit from Section~\ref{sec:vocdoni-protocol:circuits:results} to
	generate the results proof:
	\begin{align*}
		\resultsproof = {} & \hlget{\PProve}(\hlget{\resultscircuit},
		\witness = (\encresults, \epk,\\ &\merkleInclusionPf^{\encresults},
		\merkleInclusionPf^{\epk}, \{\pi^{\mathrm{dec}}_j\}_{j=1}^{n_f}), \\
		& \publicinputs = (\stateroot, \{\text{results}_j\}_{j=1}^{n_f})).
	\end{align*}
	\item Submit the results and $\resultsproof$ to the \hlget{\electionregistry} smart contract.
	The contract verifies the proof with the \hlget{\resultsverifier}, checks that the state root
	committed in the proof matches the current on-chain state root, and records the results,
	marking the process as finalized.
\end{enumerate}

\noi The election is considered finalized once the decrypted results are available on-chain. At
this stage, both the results and the complete integrity of the process are permanently recorded
and publicly auditable: anyone can recompute the on-chain commitments and check the \ac{ZK-SNARK}
proofs. In particular, the combination of the final state root, the verified decryption of the
aggregated results, and the publicly verifiable \ac{ZK-SNARK} state and results proofs guarantees
the integrity of the entire election.

%\begin{figure}%[H]
%	\centerline{\includegraphics[width=370pt,draft=false]{\figs/protocol-flow.pdf}}
%	\caption{Vocdoni voting process overview. {\color{blue}Figure is OUTDATED.}}
%	\label{fig:protocol-flow}
%\end{figure}